\section{Examples}\label{sec:examples}

The examples show the counts of Table~\ref{tab:named} at work and
connect them with the classical lists. Order three was begun in
Section~\ref{sec:ks-rational}.

\subsection{Order two: KdV--Burgers}\label{sec:kdvb}

The equation \(u_t+uu_x+\nu u_{xx}+\delta u_{xxx}=0\), \(\delta\ne0\),
has \(p=2\); the coefficient \(\nu\) may have either sign. It arises,
for instance, for waves in a liquid with gas bubbles
\citep{KudryashovSinelshchikovFourthOrder2010} and in a viscoelastic
tube \citep{KudryashovSinelshchikovChernyavsky2008}. Its exact front
solutions have been studied in several equivalent forms
\citep{Kudryashov1988,JeffreyXu1989,VliegHalford1991,Kudryashov2009Errors,Kudryashov2009KdVB}.
For \(A(\rho)=\rho+a\), polynomial division gives
\begin{equation}\label{eq:kdvb-matching}
 \remA S_A=-\frac{a(a-1)(a+1)}6,\qquad
 P_A=\rho^2+5a\rho+a^2-6a-1,\qquad
 v_A=12\bigl((a+1)Q-Q^2\bigr).
\end{equation}
Thus there are exactly \(N_2=3\) simple pairs, as
Theorem~\ref{thm:prime} states for \(2p-1=3\):
\begin{equation}\label{eq:kdvb-pairs}
 \begin{array}{c|c|c}
 a&P_A(\rho)&v_A\\ \hline
 -1&(\rho-2)(\rho-3)&-12Q^2\\
 0&\rho^2-1&12Q(1-Q)=3\operatorname{sech}^2(z/2)\\
 1&(\rho-1)(\rho+6)&12Q(2-Q).
 \end{array}
\end{equation}
The pulse \(a=0\) is the KdV soliton, the single pair of the purely
dispersive subfamily, \(F_2=1\). The two fronts belong to mixed
equations and are mirror images of each other: for a given equation
with \(\nu\ne0\) they are one and the same front, read in the two
directions (Proposition~\ref{prop:fixed}). In the original variables,
\eqref{eq:scaling} and \eqref{eq:coefficients} read
\[
 u(x,t)=w_\ast+\delta\kappa^2v_A\bigl(\kappa(x-ct)\bigr),\qquad
 \nu=\delta\kappa(5a),\qquad
 w_\ast-c=\delta\kappa^2(a^2-6a-1),
\]
so a nonzero \(\nu\) fixes the front rate to
\(\kappa=\nu/(5a\delta)\), \(a=\pm1\), in agreement with the classical
parameter restriction. The elliptic pairs at this order, both for the
KdV equation, are \eqref{eq:kdvb-elliptic}. Theorem~\ref{thm:complete}
applies to the KdV subfamily \(a_1=0\); for the mixed operators the
same conclusion cannot be expected (Remark~\ref{rem:fisher}).

\subsection{Order three: elliptic pairs, and the lattices on which they
merge}\label{sec:ks}

An elliptic solution must have zero residue at its pole. For
\(P=D^3+BD^2+\mu D+a_0\) the Laurent residue is
\(60\mu/19-15B^2/76\), so \(B^2=16\mu\); this is the curve
\(\sigma^2=16\) of \eqref{eq:ks-sigma}. With this restriction the
entire elliptic matching scheme is
\begin{equation}\label{eq:ks-elliptic-scheme}
 \begin{gathered}
 P=D^3+B D^2+\frac{B^2}{16}D+a_0,\qquad
 v=-60\wp'-15B\wp-\frac{B^3}{64}-a_0,\\
 B^4=3072g_2,\qquad
 a_0^2=\frac{13B^6}{2048}-2160g_3.
 \end{gathered}
\end{equation}
Read with the operator fixed, this is the classical solution of
\citet{Kudryashov1990}, see also \citet{Kudryashov1991}: \(w=v+a_0\)
solves the usual integrated equation
\(w'''+Bw''+\mu w'+w^2/2+K=0\) with \(K=-a_0^2/2\). For a given
equation with \(B^2=16\mu\), \(g_2\) is fixed by \(B\) and \(g_3\) is
free: a one-parameter family of elliptic waves. \citet{Hone2005} proves
that the restriction \(B^2=16\mu\) is necessary, and
\citet{Eremenko2006} completes the classification of all meromorphic KS
solutions. For the count we read \eqref{eq:ks-elliptic-scheme} the
other way: fix \(g_2,g_3\) and solve for \(B\) and \(a_0\).

On a generic lattice the equation for \(B\) has four roots and, for
each of them, the equation for \(a_0\) has two, the two backgrounds
\(\pm a_0\) of \eqref{eq:shift}: these are the \(M_3=8\) pairs. On
every lattice these two equations have eight solutions, counted with
multiplicity. This is the count of Theorem~\ref{thm:elliptic-count}:
the unknowns eliminated on the way to \eqref{eq:ks-elliptic-scheme} are
determined linearly, with nonzero rational pivots, so its last two
equations carry the full matching system, multiplicities included. On
special lattices pairs merge, and the multiplicity of a merged pair is
the number of pairs that have come together
(Section~\ref{sec:counting}). Pairs merge exactly where the Jacobian of
the two equations vanishes, and that Jacobian is a nonzero constant
times \(B^3a_0\). In terms of the invariant \(j(\Lambda)\) of
Section~\ref{sec:elliptic-counting}: if \(B=0\), then \(g_2=0\), hence
\(j=0\); if \(a_0=0\) on a lattice, elimination of \(B\) gives
\(675g_3^2=169g_2^3\), that is, \(j=-300\). These conditions are
sufficient as well. Thus:
\begin{equation}\label{eq:ks-collisions}
 \begin{array}{c|c|c}
 j(\Lambda)&\text{distinct pairs}&\text{total with multiplicity}\\ \hline
 j\notin\{0,-300\}&8&8\times1\\
 -300&6&2\times2+4\times1\\
 0&2&2\times4.
 \end{array}
\end{equation}
This is the exceptional set of Theorem~\ref{thm:elliptic-count} at
order three, and it shows the two ways in which pairs can meet. At
\(j=-300\) the four values of \(B\) remain distinct, and for two of
them the two backgrounds coincide at \(a_0=0\), in one pair of
multiplicity two; for instance, \(g_2=1/12\), \(g_3=13/1080\) gives the
double points \((B,a_0)=(\pm4,0)\). At \(j=0\) all four values of \(B\)
meet at \(B=0\), and the two roots of \(a_0^2=-2160g_3\) remain
distinct, since \(g_3\ne0\). Both pairs have the operator \(D^3+a_0\),
which is purely dissipative.

\subsection{Purely dispersive equations: Kawahara and seventh-order KdV}
\label{sec:kawahara}

\paragraph{Order four.}
For the Kawahara equation
\citep{Kawahara1972,DeminaKudryashovKawahara2010} the operator is
\begin{equation}\label{eq:kawahara-P}
 P=D^4+a_2D^2+a_0.
\end{equation}
The classification for this equation, using its first integral
\eqref{eq:energy}, is already established by
\citet{DeminaKudryashovKawahara2010}; see also \citet{YeEtAl2022}. By
Theorem~\ref{thm:symmetric}, the pairs with an even operator are those
with \(\iota A=A\), that is, \(A=\rho^3+\alpha_2\rho\). All of them
are pulses: the profile is even about its pole, \(v(Q)=v(1-Q)\). The
quotient and remainder are
\begin{equation}\label{eq:kawahara-exponential}
 \begin{split}
 P_A&=\rho^4+13\alpha_2\rho^2+\frac{31\alpha_2^2-70\alpha_2+7}{3},\\
 v_A&=280(1+\alpha_2)Q(1-Q)-1680Q^2(1-Q)^2,\\
 \remA S_A
   &=-\frac{\rho}{420}(\alpha_2+1)(31\alpha_2^2-31\alpha_2+10).
 \end{split}
\end{equation}
The three distinct roots, \(\alpha_2=-1\) and
\(\alpha_2=1/2\pm3i/(2\sqrt{31})\), give \(F_4=\binom31=3\) simple
pairs. The cubic already appears in \citet[Eq.~(47)]{Kudryashov1989},
with \(y=-1/\alpha_2\), where one real and two complex conjugate roots
are noted. The three values are the simply periodic parameter choices
of \citet[Section~3]{DeminaKudryashovKawahara2010}, whose parameter in
Eq.~(31) equals \(1/\alpha_2\). At \(\alpha_2=-1\),
\begin{equation}\label{eq:kawahara-soliton}
 v=-105\operatorname{sech}^4(z/2),\qquad
 P=(D^2-4)(D^2-9).
\end{equation}
This is the familiar solitary wave of the Kawahara equation; the other
two pairs are complex at unit real rate. Of all \(N_4=35\) pairs at
this order, \(21\) have real coefficients, and all \(21\) have rational
coefficients: five pulses and sixteen fronts, eleven wave shapes up to
reflection. They are the one-pole waves of the mixed fifth-order
equations, such as the Benney--Lin equation and the equation of
\citet{KudryashovSinelshchikov2011} for a viscoelastic tube, and the
repository dictionary lists them; solitary waves of this family are
found in \citet{KudryashovDemina2009}.

Substitution shows that the elliptic pairs of \eqref{eq:kawahara-P} are
given by
\begin{align}
 v&=-1680\wp^2-\frac{280a_2}{13}\wp
       +168g_2+\frac{31a_2^2}{507}-a_0,
       \label{eq:kawahara-elliptic}\\
 0&=31a_2^3-42588a_2g_2-4745520g_3,
       \label{eq:kawahara-cubic}\\
 a_0^2&=23184g_2^2+\frac{7980}{169}a_2^2g_2
                       +\frac{101520}{13}a_2g_3.
       \label{eq:kawahara-background}
\end{align}
These are the Weierstrass profiles and parameter relations obtained
from Laurent series by
\citet[Section~3]{DeminaKudryashovKawahara2010}; in their notation
\(a_2=-\beta\), \(a_0^2=C_0^2+12C_1\) and \(v=-6(w-w_\ast)\), where
their \(w_\ast\) is a constant equilibrium and \(a_0=C_0-6w_\ast\).
The first condition is cubic in \(a_2\), and the second is monic
quadratic in \(a_0\), so there are \(3\times2=6\) solutions, counted
with multiplicity, on every lattice: three values of \(a_2\), each with
its two backgrounds. This is the number \(2F_4\) of
Corollary~\ref{cor:elliptic-symmetric}. At \((g_2,g_3)=(1,0)\) the
cubic has three distinct roots and the right side of
\eqref{eq:kawahara-background} is nonzero at all of them, so the six
pairs are simple on a nonempty open set of lattices.
The unrestricted elliptic count at this order is \(M_4=30\); elliptic
formulas with nonzero odd derivative coefficients in the full
fifth-order evolution equation are given by
\citet{KudryashovSinelshchikov2012}.

\paragraph{Order six.}
The seventh-order KdV equation has the operator
\(P=D^6+a_4D^4+a_2D^2+a_0\); its meromorphic travelling waves are
studied by \citet{Dang2023Even}.
Here \(A=\rho^5+\alpha_2\rho^3+\alpha_4\rho\) and \(s_6=5544\). The
remainder of \(S_A\) has two coefficients, at \(\rho^3\) and at
\(\rho\), which are polynomials in \(\alpha_2,\alpha_4\) of weights
eight and ten. Eliminating \(\alpha_4\), which is determined by
\(\alpha_2\) at every solution, gives
\begin{equation}\label{eq:p6-eliminant}
 (\alpha_2+5)(\alpha_2+10)(49\alpha_2^2-5\alpha_2+10)\,
 \varphi_6(\alpha_2)=0,
\end{equation}
where \(\varphi_6\) is an irreducible sextic without real roots; the
ancillary script prints it and the two coefficients. The ten roots are
distinct, in agreement with \(F_6=\binom52=10\), and all ten pairs are
pulses. Two pairs are real, the solitary waves of
\citet{RyabovSinelshchikovKochanov2011}:
\begin{equation}\label{eq:p6-real}
 \begin{array}{c|c|c}
 A&P_A&v_A\\ \hline
 \rho(\rho^2-1)(\rho^2-4)&(\rho^2-9)(\rho^2-16)(\rho^2-25)
   &10395\operatorname{sech}^6(z/2)\\
 \rho(\rho^2-1)(\rho^2-9)&(\rho^2-4)(\rho^2-25)(\rho^2-71)
   &10395\operatorname{sech}^4(z/2)\bigl(1+\operatorname{sech}^2(z/2)\bigr).
 \end{array}
\end{equation}
An exact computation on a sample lattice gives \(2F_6=20\) for the
elliptic pairs of this family.

\subsection{Purely dissipative equations of order five}\label{sec:p5}

At \(p=5\) the purely dissipative operators are
\(P=D^5+a_3D^3+a_1D+a_0\), and \(A=\rho^4+\alpha_2\rho^2+\alpha_4\).
All \(126\) pairs at \(p=5\) are simple
(Appendix~\ref{app:verification}). By
Proposition~\ref{prop:general}(iii) there are therefore exactly
\(F_5=6\) distinct pairs here, and all six are fronts,
\(A(0)=\alpha_4\ne0\); for the four real ones this is also
Theorem~\ref{thm:symmetric}(iii). They are the three pairs in the
first line of
\begin{equation}\label{eq:p5-symmetric}
 \begin{gathered}
 A=(\rho^2-1)(\rho^2-4),\qquad
 A=(\rho^2-1)(\rho^2-9),\qquad
 A=(\rho^2-1)\Bigl(\rho^2+\frac15\Bigr),\\
 1105\alpha_2^3-2969\alpha_2^2+5570\alpha_2-4000=0,\qquad
 \alpha_4=\frac{5\alpha_2^3+21\alpha_2-10}{79\alpha_2+210},
 \end{gathered}
\end{equation}
and the three pairs given by the second line; the cubic has one real
root. The six pairs are the six solitary waves of
\citet{KudryashovDemina2007}, in whose notation \(w=1/(4a_3)\). The
four with real coefficients are, after the scales are restored, the
solutions~(1.4) of \citet{KudryashovMigita2007}; see also
\citet{RyabovSinelshchikovKochanov2011}. For instance,
\(A=(\rho^2-1)(\rho^2-4)\) gives
\(P_A=\rho^5-55\rho^3+1254\rho-2520\) and
\(v_A=5040\,Q^3(10-15Q+6Q^2)\), a front from \(0\) to
\(5040=-2P_A(0)\).
